% !TeX root = ../kaehler-einstein-degeneration.tex
% Current content: 2026-11-23, v1.3.
\section{全深部三重点領域と complete simplex KE model}
\label{coarse:sec:full-triple-comparison}

前節の product barrier は，一つの深度が他の二つより大きい領域で
明示的な cusp product への漸近を与えた．本節では product に代えて
logarithmic simplex 全体の complete toric KE metric を用いる．
この模型の境界を固定量だけ内外へ平行移動することにより，balanced な点を含む
全深部三重点領域で actual metric を比較する．
外部存在入力には，有界な\emph{実}単体上の Cheng--Yau の凸解と
実 Hessian metric の完備性を用いる
\cite[Theorem~2.1]{ZhangNegativeKE}．
複素 tube の完備性，一意性，尺度変更と以下の一様評価は，この入力から本文で導く．
有界な実底空間と，虚方向には非有界な複素 tube を区別する．
同論文の退化族に関する Theorem~2.4 や，その全空間の
$\mathbb Q$-factorial 仮定は用いない．
まず exact triple 領域の固定 face 上に一様な density 下界を持つ負 KE metrics に対し，
局所安定性を証明する．各固定 parameter で radial closure と全 phase tori が compact であり，
metric がその近傍で滑らかなことを仮定するが，parameter 一様な境界上界は要しない．
その後，\cref{ass:Gd} を満たす固定次数 $d\ge5$ の全平面退化へ，
\cref{coarse:cor:degree-d-volume-lower} を用いて適用する．
Actual estimates の定数は固定した $d$ と退化族に依存してよく，次数に一様な評価は主張しない．
以下の比較法と本族への適用について，新規性の包括的確認は未完了である．

\subsection{Complete simplex model とその一様幾何}

$L>0$ に対して
\[
 \Delta_L=\{(\ell_0,\ell_1,\ell_2):
              \ell_i>0,\ \ell_0+\ell_1+\ell_2=L\}
\]
と置く．独立な実座標は $(\ell_1,\ell_2)$ とする．
対応する tube domain では $w_0+w_1+w_2=L$，$\Re w_i=\ell_i>0$ と書く．
三重点の angular universal cover 上では，imaginary constants を調整した
$w_i=-2\log\zeta_i+\text{imaginary constant}$ を使える．
$\eta_i=d\log\zeta_i=-dw_i/2$ とし，前節までと同じ
\[
 dV_{\log}
  =(\sqrt{-1}\eta_1\wedge\bar\eta_1)
    \wedge(\sqrt{-1}\eta_2\wedge\bar\eta_2),
 \qquad
 \mathcal V_L=\frac{8L^2}{\ell_0^2\ell_1^2\ell_2^2}\,dV_{\log}
\]
を用いる．以下では $\omega_L^{\mathrm{cusp}}$ を
\eqref{coarse:eq:triple-cusp-reference} の constrained cusp form とする．

\begin{lemma}[Complete simplex KE model]
\label{coarse:lem:simplex-model}
$\Delta_L$ 上に一意な滑らかな strictly convex function $\phi_L$ が存在して
\begin{equation}\label{coarse:eq:simplex-real-MA}
 \det D^2\phi_L=e^{2\phi_L},\qquad
 \phi_L|_{\partial\Delta_L}=+\infty
\end{equation}
を満たす．Tube 上の form
\begin{equation}\label{coarse:eq:simplex-KE-form}
 k_L=2\ddc\phi_L
    =\frac{\sqrt{-1}}2\sum_{j,k=1}^2
       (\phi_L)_{jk}\,dw_j\wedge d\bar w_k
\end{equation}
は complete negative KE metric であり，imaginary translations に不変なので
全 phase tori の quotient に降りる．
\[
 \psi_L=\log L-\sum_{i=0}^2\log\ell_i
\]
と置くと，次の評価が成立する：
\begin{align}
 \psi_L-\tfrac12\log3&\le\phi_L\le\psi_L,
        \label{coarse:eq:simplex-potential-bounds}\\
 \tfrac13\mathcal V_L&\le k_L^2\le\mathcal V_L,
        \label{coarse:eq:simplex-density-bounds}\\
 C^{-1}\omega_L^{\mathrm{cusp}}&\le k_L
                       \le C\omega_L^{\mathrm{cusp}}.
        \label{coarse:eq:simplex-model-tensor}
\end{align}
定数 $C$ は $L$ と点の深度比に依存しない．
最大深度を従属座標に選ぶ固定相対窓の normalized holomorphic coordinates 上では，
$k_L$ は一様な ellipticity と全固定階の $C^q$ bounds を持つ．
この一様性は angular universal cover の固定半径座標でのものであり，
固定周期の phase quotient の injectivity radius の下界を主張しない．
\end{lemma}

\begin{proof}
\emph{用いる外部存在結果．}
Zhang\cite[Theorem~2.1]{ZhangNegativeKE} が Cheng--Yau の実凸領域定理から
取り出す入力は，次のものである．有界な実単体 $\Delta_1$ 上に，
$\det D^2\phi_1=e^{2\phi_1}$，$\phi_1\to+\infty$ at $\partial\Delta_1$ を満たす
滑らかな strictly convex 解が少なくとも一つ存在し，
\[
 g_{B,1}=\sum_{j,k=1}^2(\phi_1)_{jk}\,d\ell_j\,d\ell_k
\]
は complete である．多面体の境界を滑らかな strongly convex 境界とみなすことも，
非有界な tube に有界 domain の存在定理を直接適用することも行わない．
この入力から先は次の直接証明を用いる．

\emph{尺度変更と一意性．}
一般の $L>0$ には
\begin{equation}\label{coarse:eq:simplex-scaling}
 \phi_L(\ell)=\phi_1(\ell/L)-2\log L
\end{equation}
を用いる．$\phi_L$ は実部のみに依存するため，translation invariance は直接従う．
二実次元の Hessian determinant は $L^{-4}$ 倍になるので，
\eqref{coarse:eq:simplex-real-MA} と境界 blow-up が保たれる．
$g_{B,L}$ は $\ell\mapsto\ell/L$ による $g_{B,1}$ の pullback であり complete である．
独立な二つの深度の選び直しは絶対値 $1$ の affine Jacobian を持つため，
方程式とこの規格化はその選択に依存しない．

一意性は simplex の内側 homothety だけから従う．
$u,v$ を \eqref{coarse:eq:simplex-real-MA} を満たす二解とし，
$0<\varepsilon<L/3$，$a=1-3\varepsilon/L$ に対して
\[
 v_\varepsilon^{\mathrm{in}}(\ell)
 =v\bigl((\ell-\varepsilon\mathbf1)/a\bigr)-2\log a
 \quad(\ell_i>\varepsilon)
\]
と置く．引数の三深度の和は $L$ であり，
$\det D^2v_\varepsilon^{\mathrm{in}}=e^{2v_\varepsilon^{\mathrm{in}}}$ である．
内側単体の境界では $v_\varepsilon^{\mathrm{in}}\to+\infty$，
その閉包で $u$ は有限である．$u-v_\varepsilon^{\mathrm{in}}$ の正の最大点では
Hessian の大小と Monge--Amp\`ere 方程式が矛盾するので
$u\le v_\varepsilon^{\mathrm{in}}$ である．$\varepsilon\downarrow0$ から $u\le v$，
両者を交換して $u=v$ を得る．境界発散率や二つ目の解の完備性を追加仮定しない．

\emph{複素 tube と phase quotient の完備性．}
$w_j=\ell_j+\sqrt{-1}y_j$ と書くと，$k_L$ の実計量は
\begin{equation}\label{coarse:eq:simplex-real-block-metric}
 g_L^{\mathrm{sf}}
 =\sum_{j,k=1}^2(\phi_L)_{jk}
       (d\ell_j\,d\ell_k+dy_j\,dy_k).
\end{equation}
従って実底への射影は長さを増やさない．有限長の曲線の実射影は，
$g_{B,L}$ の完備性により $\Delta_L$ の内点へ収束する．
その末尾の実射影を含む compact 集合では $D^2\phi_L$ に正の下界があるので，
虚座標の変動も Euclidean 長で有限であり，有限点へ収束する．
よって tube に有限長の発散曲線はなく，$g_L^{\mathrm{sf}}$ は complete である．
また
\[
 \Ric(k_L)=-\ddc\log\det D^2\phi_L=-2\ddc\phi_L=-k_L.
\]
虚平行移動の格子は free かつ properly discontinuous に等長的に働くため，
phase quotient も complete である．Actual 座標 $w_i=-2\log\zeta_i+\mathrm{const.}$
では各独立虚座標の周期は $4\pi$ である．
写像 $w\mapsto w/L$ は tube の universal covers の間では等長だが，
この格子を周期 $4\pi/L$ の格子へ送る．従って固定周期 $4\pi$ の quotient 同士を
同じ尺度で等長と主張しているわけではない．

\emph{Potential と density の障壁．}
直接計算により
\begin{equation}\label{coarse:eq:simplex-barrier-determinant}
 \frac{\det D^2\psi_L}{e^{2\psi_L}}
   =\frac{\ell_0^2+\ell_1^2+\ell_2^2}{L^2}
   \in[1/3,1].
\end{equation}
$0<\varepsilon<L/3$ として，内側 simplex $\ell_i>\varepsilon$ 上で
\[
 \psi^{\mathrm{in}}_\varepsilon(\ell)
 =\psi_{L-3\varepsilon}(\ell-\varepsilon\mathbf1)
\]
を考える．その境界では $\psi^{\mathrm{in}}_\varepsilon\to+\infty$ である一方，
$\phi_L$ はその閉包で有限である．もし
$\phi_L-\psi^{\mathrm{in}}_\varepsilon$ が正の最大値を持てば，その内点で
\[
 e^{2\phi_L}=\det D^2\phi_L
 \le\det D^2\psi^{\mathrm{in}}_\varepsilon
 \le e^{2\psi^{\mathrm{in}}_\varepsilon}
\]
となり矛盾する．従って $\phi_L\le\psi^{\mathrm{in}}_\varepsilon$ である．

逆に，元の $\Delta_L$ 上で
\[
 \psi^{\mathrm{out}}_\varepsilon(\ell)
 =\psi_{L+3\varepsilon}(\ell+\varepsilon\mathbf1)
   -\tfrac12\log3
\]
を考える．この関数は $\overline{\Delta_L}$ 上で有限であり，
$\phi_L\to+\infty$ on $\partial\Delta_L$ である．
$\psi^{\mathrm{out}}_\varepsilon-\phi_L$ の正の内部最大点が存在すれば
\[
 e^{2\psi^{\mathrm{out}}_\varepsilon}
 \le\det D^2\psi^{\mathrm{out}}_\varepsilon
 \le\det D^2\phi_L=e^{2\phi_L}
\]
となり矛盾する．$\varepsilon\downarrow0$ とすることにより
\eqref{coarse:eq:simplex-potential-bounds} を得る．

Volume の係数を確認する．ここで
$dV_w=(\sqrt{-1}\,dw_1\wedge d\bar w_1)
\wedge(\sqrt{-1}\,dw_2\wedge d\bar w_2)$ と定義すれば，
$dV_{\log}=\tfrac1{16}dV_w$ であり，
\eqref{coarse:eq:simplex-KE-form} から
\begin{equation}\label{coarse:eq:simplex-volume-normalization}
 k_L^2=8\det D^2\phi_L\,dV_{\log}
      =8e^{2\phi_L}\,dV_{\log}.
\end{equation}
従って potential bounds は \eqref{coarse:eq:simplex-density-bounds} を与える．

Tensor の下比較には complete model での Schwarz lemma を用いる．
三つの正則写像
\[
 f_i:(\Delta_L+\sqrt{-1}\RR^2,k_L)
      \longrightarrow\{\Re w>0\},\qquad f_i(w)=w_i
\]
を考え，target には曲率 $-1$ の half-plane metric を入れる．
Source は complete で $\Ric(k_L)=-k_L$ なので，
Yau\cite[Theorem~2]{YauSchwarzThm2} により
$f_i^*\omega_{\mathbb H}\le C_0 k_L$ である．
三つを加えると $\omega_L^{\mathrm{cusp}}\le3C_0 k_L$ となる．
これは complete model 上の議論であり，actual metric の completeness や
局所境界までの actual distance を用いていない．

逆側について，\eqref{coarse:eq:triple-cusp-exact-volume-ratio} と
\eqref{coarse:eq:simplex-density-bounds} より
\[
 \frac{k_L^2}{(\omega_L^{\mathrm{cusp}})^2}\le3.
\]
複素次元が $2$ なので，先に得た各固有値の正の下界と determinant の上界から，
各固有値の一様上界を得る．これで \eqref{coarse:eq:simplex-model-tensor} が成立する．

\emph{全深度比に一様な normalized geometry．}
任意の中心 $p$ で最大深度を $c=\ell_i(p)$，残る二深度を
$a=\ell_j(p)$，$b=\ell_k(p)$ とし，$\delta=1/8$ を固定する．
\[
 \xi_j=\frac{w_j-w_j(p)}a,\qquad
 \xi_k=\frac{w_k-w_k(p)}b,\qquad x=\Re\xi_j,\quad y=\Re\xi_k
\]
と置くと，$|x|,|y|<\delta$ の固定 strip product 上で
\[
 \ell_j=a(1+x),\quad\ell_k=b(1+y),\quad
 \ell_i=c-ax-by\in[(1-2\delta)c,(1+2\delta)c]
\]
となる．これは任意の正の深度に対して tube 内に含まれ，
actual 三重点 chart の大きさや任意半径の球の包含を要求しない．
この座標で constrained cusp metric は
\[
 \frac{|d\xi_j|^2}{(1+x)^2}
 +\frac{|d\xi_k|^2}{(1+y)^2}
 +\frac{|a\,d\xi_j+b\,d\xi_k|^2}{(c-ax-by)^2}
\]
である．$a,b\le c$ なので，これは固定 Euclidean metric と一様に比較できる．
従って既に証明した \eqref{coarse:eq:simplex-model-tensor} により，
$k_L$ の pullback の coefficient matrix も一様に正で有界である．
Potential の規格化も明示すると
\begin{equation}\label{coarse:eq:simplex-normalized-potential}
 \Phi_p(x,y)=\phi_L\bigl(\ell(x,y)\bigr)+\log(ab)
\end{equation}
は $\det D^2\Phi_p=e^{2\Phi_p}$ を満たし，
\eqref{coarse:eq:simplex-potential-bounds} と $1\le L/c\le3$ から
固定 strip product の実底上で一様に有界である．

各 strip の complete curvature $-1$ metric の積を固定した比較背景とする．
内側半 strip 上ではその全階係数は一様に有界で，Euclidean metric と一様同値である．
従って $k_L$ との tensor 二側比較を既知の入力として，さらに小さい固定座標球で
\cref{coarse:lem:nested-metric} を適用し，全固定階の bounds を得る．
ここで用いる rectangles は内部正則性の小さい窓だけであり，
単体全体や任意半径の球を一つの積領域へ置き換えてはいない．
Tensor 比較は完備模型への Schwarz 評価と determinant 上界から先に得ており，
bounded potential 又は bounded density 単独から導いたものではない．
\end{proof}

\subsection{平行移動した simplex の density 比較}

\begin{lemma}[Translated-simplex density estimate]
\label{coarse:lem:simplex-shift}
Logarithmic frame に関する model density を
\begin{equation}\label{coarse:eq:simplex-log-density}
 \rho_L(\ell)=\log\frac{k_L^2}{dV_{\log}}
             =\log8+2\phi_L(\ell)
\end{equation}
とする．$c>0$, $L\ge6c$, $m:=\min_i\ell_i\ge2c$ ならば
\begin{equation}\label{coarse:eq:simplex-shift-bound}
 \left|\rho_{L\pm3c}(\ell\pm c\mathbf1)-\rho_L(\ell)\right|
 \le C\frac{c}{m}.
\end{equation}
定数は $L,c,\ell$ に依存しない．
\end{lemma}

\begin{proof}
\eqref{coarse:eq:simplex-scaling} により
\[
 \rho_L(\ell)=\rho_1(\ell/L)-4\log L.
\]
$\rho_1$ は凸であり，\cref{coarse:lem:simplex-model} より
\[
 -\log3\le
 \rho_1(x)-\left(\log8-2\sum_i\log x_i\right)\le0.
\]
各 $x\in\Delta_1$ で最大 $x_i$ を従属に選び，独立な二変数 $x_j,x_k$ の
相対幅 $1/8$ の rectangle を取る．この rectangle 全体で三つの coordinates は
中心値と固定比の間にあるため，$-2\sum\log x_i$ の oscillation は一様有界である．
従って $\rho_1$ の oscillation も一様有界である．
凸関数の一次差分による gradient estimate から
\[
 |\partial_{x_j}\rho_1(x)|\le C/x_j,
 \qquad |\partial_{x_k}\rho_1(x)|\le C/x_k
\]
を得る．ここで微分は $x_i=1-x_j-x_k$ として取る．
従って任意の tangent vector $v$，$\sum_i v_i=0$ に対し
\begin{equation}\label{coarse:eq:simplex-log-gradient}
 |d\rho_1(x)(v)|\le C\sum_{i=0}^2\frac{|v_i|}{x_i}.
\end{equation}

$x=\ell/L$ とし，$|s|\le c/L$ の区間で
\[
 q_s=\frac{x+s\mathbf1}{1+3s}
\]
と置く．仮定より $|s|\le\frac12\min x_i$，$1+3s\ge1/2$ であり，
$q_{s,i}\ge x_i/2$ と $|\dot q_{s,i}|\le8$ が成立する．
\eqref{coarse:eq:simplex-log-gradient} により
\[
 \left|\frac{d}{ds}\rho_1(q_s)\right|
 \le\frac{C}{\min x_i}.
\]
$s=0$ から $s=\pm c/L$ まで積分し，
$|\log(1\pm3c/L)|\le Cc/L$ も用いると
\[
 \begin{split}
 &\left|\rho_{L\pm3c}(\ell\pm c\mathbf1)-\rho_L(\ell)\right|\\
 &\quad=
 \left|\rho_1(q_{\pm c/L})-\rho_1(x)
             -4\log(1\pm3c/L)\right|
 \le Cc/(L\min x_i)=Cc/m.
 \end{split}
\]
\end{proof}

\subsection{固定 face 下界からの局所安定性}

\begin{theorem}[Exact triple 領域の境界密度からの局所安定性]
\label{coarse:thm:boundary-to-triple}
$m_0,c_*>0$ を固定し，$0<|t|<t_0$ に対し
\[
 M_t=\{(\zeta_0,\zeta_1,\zeta_2)\in(\CC^*)^3:
                  \zeta_0\zeta_1\zeta_2=t\},\qquad
 \ell_i=-\log|\zeta_i|^2,\qquad L_t=-\log|t|^2
\]
とする．$L_t>3m_0$ とし，全 angular phases を含む領域
\[
 D_{t,m_0}=\{\zeta\in M_t:\ell_i>m_0\ (i=0,1,2)\}
\]
を考える．各 $t$ について，$\omega_t$ は $M_t$ 内の
$\overline{D_{t,m_0}}$ の近傍で滑らかな Kähler metric であり，
$\Ric(\omega_t)=-\omega_t$ とする．さらに全ての境界 face と angular phases 上で
\begin{equation}\label{coarse:eq:local-triple-face-floor}
 \omega_t^2\ge c_*\mathcal V_{L_t}
 \qquad\text{on }\partial D_{t,m_0}
\end{equation}
を仮定する．

このとき，$m_0,c_*$ のみに依存する $m_1>m_0$ と $L_*>0$，
および各整数 $q\ge0$ に対する $C_q$ が存在し，
$L_t\ge L_*$，$m_1\le m<L_t/3$ ならば
\begin{align}
 \sup_{\min_i\ell_i\ge m}
   \left|\log\frac{\omega_t^2}{k_{L_t}^2}\right|
  &\le C_0/m,
       \label{coarse:eq:local-triple-small-volume}\\
 \|g_t-g_{L_t}^{\mathrm{sf}}\|_{C^q(\{\min_i\ell_i\ge m\},g_{L_t}^{\mathrm{sf}})}
  &\le C_q/m.
       \label{coarse:eq:local-triple-small-tensor}
\end{align}
ここで $g_t,g_{L_t}^{\mathrm{sf}}$ は $\omega_t,k_{L_t}$ の実計量であり，
$C^q$ norm は最大深度を従属にする normalized holomorphic covering charts，
または同値な model covariant norm で計算する．定数は全深度比・全 angular phases に一様であり，
$C_q$ は $m_0,c_*,q$ のみに依存する．
各固定 $t$ の滑らかな閉包への延長以外に $t$ に関する正則性は仮定せず，
境界上の density 上界を $t$ に一様に取る必要もない．
\end{theorem}

\begin{proof}
以下 $L=L_t$ と書く．最大値原理は $D_{t,m_0}$ 自身，すなわち全 angular phases を含む
compact radial closure 上で用いる．Angular universal cover は最後の局所正則性にのみ用いる．

まず下 barrier を作る．$c=M m_0$ とし，$M\ge1$ を十分大きい固定定数に取る．
平行移動は $w_i\mapsto w_i+c$ という正則写像なので，
expanded simplex model の引戻し
\[
 k_L^{\mathrm{out}}(\ell)=k_{L+3c}(\ell+c\mathbf1)
\]
は $D_{t,m_0}$ の閉包の近傍で滑らかな負 KE metric である．
\Cref{coarse:lem:simplex-model} により
\begin{equation}\label{coarse:eq:expanded-simplex-face-ratio}
 \frac{(k_L^{\mathrm{out}})^2}{\mathcal V_L}
 \le\left(1+\frac{3c}{L}\right)^2
       \prod_{i=0}^2\left(\frac{\ell_i}{\ell_i+c}\right)^2.
\end{equation}
境界では少なくとも一つの $\ell_i$ が $m_0$ に等しく，対応する因子は
$(1+M)^{-2}$ である．他の二因子は $1$ 以下である．
$L\ge6c$ とし，$M$ を $\tfrac94(1+M)^{-2}\le c_*$ となるように選べば
\[
 (k_L^{\mathrm{out}})^2\le c_*\mathcal V_L\le\omega_t^2
 \qquad\text{on }\partial D_{t,m_0}.
\]
両 metric の Ricci curvature は $-1$ であるから，
$\log(\omega_t^2/(k_L^{\mathrm{out}})^2)$ の負の内部最小点では
$\omega_t-k_L^{\mathrm{out}}\ge0$ となり，determinant の符号に矛盾する．
従ってこの volume 下比較は $D_{t,m_0}$ 全体で成り立つ．

上 barrier には inner simplex 自身の complete model
\[
 k_L^{\mathrm{in}}(\ell)=k_{L-3m_0}(\ell-m_0\mathbf1)
\]
を用いる．\eqref{coarse:eq:simplex-density-bounds} よりその density は
$\partial D_{t,m_0}$ で無限大へ発散する．
各固定 $t$ で $\omega_t$ の logarithmic density は閉包上で有限であるから，
$\log(\omega_t^2/(k_L^{\mathrm{in}})^2)$ は境界で $-\infty$ へ向かう．
この有限上界は $t$ に依存してよく，以下の比較定数には現れない．
正の内部最大点における同じ determinant 比較から上比較を得る．
以上より
\begin{equation}\label{coarse:eq:full-simplex-volume-sandwich}
 (k_L^{\mathrm{out}})^2\le\omega_t^2\le(k_L^{\mathrm{in}})^2
 \qquad\text{on }D_{t,m_0}.
\end{equation}
この最大値原理は compact phase quotient 上で適用しており，
actual distance に関する境界逃避や tensor 比較を前提としていない．

$m\ge2c$ とし，\cref{coarse:lem:simplex-shift} を下側の shift $+c$ と
上側の shift $-m_0$ に用いる．$m_0,c$ は固定されているので
\eqref{coarse:eq:full-simplex-volume-sandwich} は
\eqref{coarse:eq:local-triple-small-volume} を与える．

最後に，$p$ が $\min_i\ell_i(p)\ge m$ を満たすとする．
最大深度を従属にする固定相対窓では，全ての深度は $3m/4$ 以上であり，
同じ volume squeeze による
\[
 v_t=\log(\omega_t^2/k_L^2),\qquad \|v_t\|_{C^0}\le C/m
\]
がその窓全体で成立する．
\Cref{coarse:lem:simplex-model} により背景 $k_L$ はその normalized charts で
一様な ellipticity と全固定階の bounds を持つ．
両 metric は負 KE なので
\[
 \ddc v_t=\omega_t-k_L,\qquad
 (k_L+\ddc v_t)^2=e^{v_t}k_L^2.
\]
$m_1\ge4c$ をさらに十分大きく固定すれば，
\cref{coarse:lem:small-volume-tensor} の Savin small-perturbation estimate が
全ての窓に一様に適用できる．同補題の線形化後の Schauder estimates により
\eqref{coarse:eq:local-triple-small-tensor} を得る．
\end{proof}

\subsection{固定次数の全平面退化への適用}

\Cref{ass:Gd} を満たす固定次数 $d\ge5$ の全平面退化 $X_t^{(d)}$ を取り，
$\Ric(\omega_t)=-\omega_t$ と規格化した actual KE metric を考える．
Source を含まない三重点 chart の exact coordinates
$\zeta_0\zeta_1\zeta_2=t$ に戻り，$L=L_t$ とする．
以下の固定 $m_0$ は，$\min_i\ell_i\ge m_0$ の閉領域が chart の内部にあるよう，
十分大きく取る．三つの active depths では
$\alpha_{i,t}=A_0+\min\{\ell_i,L/2\}+O(1)\asymp\ell_i$ であり，
残る $d-3$ 個の depths と canonical residue frame の係数は正の定数の間にある．
従って $\Omega_{t,d}^{\mathrm{low}}\asymp\mathcal V_L$ であり，
\cref{coarse:cor:degree-d-volume-lower} により，ある一様定数 $c_*>0$ に対して
\begin{equation}\label{coarse:eq:actual-simplex-density-floor}
 \omega_t^2\ge c_*\mathcal V_L
 \qquad\text{on }\{\min_i\ell_i\ge m_0\}
\end{equation}
である．定数は固定した $d$ と族に依存し，有限個の三重点について一様に取れる．

\begin{theorem}[全深部三重点領域の actual semi-flat tensor 漸近]
\label{coarse:thm:full-triple-tensor}
\Cref{ass:Gd} を満たす固定次数 $d\ge5$ の全平面退化の三重点で，
ある固定 $m_1>m_0$ と各 $q\ge0$ に対する
$C_q$ が存在し，十分小さい $t\ne0$ と $m_1\le m<L_t/3$ に対し
\begin{align}
 \sup_{\min_i\ell_i\ge m}
   \left|\log\frac{\omega_t^2}{k_{L_t}^2}\right|
  &\le C/m,
       \label{coarse:eq:full-triple-small-volume}\\
 \|g_t-g_{L_t}^{\mathrm{sf}}\|_{C^q(\{\min_i\ell_i\ge m\},g_{L_t}^{\mathrm{sf}})}
  &\le C_q/m.
       \label{coarse:eq:full-triple-small-tensor}
\end{align}
ここで $g_{L_t}^{\mathrm{sf}}$ は $k_{L_t}$ の Riemannian metric である．
Norm は最大深度を従属にする normalized holomorphic covering charts で計算し，
同値な model covariant norms でも評価できる．
評価は全ての深度比と全 angular phases について一様である．
また固定した族の全三重点に一様であり，定数は $d$ と退化族に依存してよい．
特に $m_t\to\infty$ としても，balanced region を含む
$\min_i\ell_i\ge m_t$ の全体で actual/model comparison が成立する．
\end{theorem}

\begin{proof}
十分大きい固定 $m_0$ を取れば，各固定 $t\ne0$ の全 phase を含む
$\overline{D_{t,m_0}}$ は exact triple chart に compact に含まれる．
Actual KE metric はその近傍で滑らかであり，
\eqref{coarse:eq:actual-simplex-density-floor} は
\cref{coarse:thm:boundary-to-triple} の境界入力
\eqref{coarse:eq:local-triple-face-floor} を与える．同定理を各 triple chart へ適用する．
固定族の三重点は有限個なので，$m_1,L_*,C_q$ を共通に取れる．
\end{proof}

\begin{remark}[Actual 比較の最小境界入力]
\label{coarse:rem:full-triple-face-input}
\Cref{coarse:thm:boundary-to-triple} は exact triple 領域だけの定理であり，
大域的な compactification や全平面退化を仮定しない．
\Cref{coarse:cor:degree-d-volume-lower} は本族への適用で
\eqref{coarse:eq:local-triple-face-floor} を供給する一つの前段入力である．
Interior の density 下界，actual tensor 比較，actual distance に関する境界逃避，
または actual 領域の完備性は別途仮定しない．
一方，最大値原理には全 phase の compact quotient と各固定 $t$ の閉包への滑らかな延長を使う．
非周期な無限 cover 上の局所 metrics や，一部の角度だけの領域へ，
同じ境界仮定だけで定理を適用することはできない．
本稿の angular universal cover 上の評価は quotient 上の metrics の引戻しである．

小さい potential は，shift 評価で小さくなった
$v_t=\log(\omega_t^2/k_{L_t}^2)$ 自身である．
同じ負 KE 規格化による Ricci 恒等式
$\ddc v_t=\omega_t-k_{L_t}$ が相対 potential を正確に与えるので，
有界な density だけから potential の振幅を小さく選び直す議論は含まれない．
\end{remark}

\begin{remark}[この比較が解決する範囲]
\label{coarse:rem:full-triple-scope}
\Cref{coarse:lem:simplex-model,coarse:thm:full-triple-tensor} を合わせると，
全深部三重点領域で actual metric と constrained cusp metric の一様二側比較が得られる．
従ってこの領域での短い torus averaging--Moser 構成には，
別途の actual tensor-comparison 仮定を要しない．
Sharp な $O(m^{-1})$ 比較の相手は complete semi-flat metric $k_{L_t}$ である．
Balanced な点で constrained cusp form 自体を KE metric と同一視する主張ではない．

本節は source cores，source-arm の全体，異なる三重点 chart の相対接着，
領域外を通る shortcut，または大域 GH 極限を扱っていない．
したがって exterior 全体の tensor 比較や，大域条件
$\mathrm{(CP)}_d+\mathrm{(PBE)}_d$ の検証とは区別する．
Fixed toric model 自体の定数は次数に依存しない．一方，actual family について用いる
\eqref{coarse:eq:actual-simplex-density-floor} の定数，chart の選択，
およびそこから得る actual comparison の定数には固定した $d$ と族への依存を許す．
次数を動かした一様な estimate は本節の結論に含まれない．
\end{remark}
